
This work addresses the problem of deep learning researchers wasting time testing hypotheses that prove infeasible due to missing datasets, incompatible benchmarks, or unavailable automated metrics. The constraint-satisfiability verification system checks for (Dataset, Benchmark, Metric) triple existence in a structured knowledge base and flags cross-domain confound patterns before experiments begin.

Experimental ground truth validation---testing 20 system-classified ''testable'' hypotheses and measuring post-hoc p-values---demonstrates 90\% experimental success rate (binomial p = 0.0002 vs. 50\% random baseline), exceeding the 75\% prediction threshold by 15 percentage points. This non-circular evaluation approach validates predictions against actual experimental outcomes rather than circular expert consensus. Cross-domain confound pattern transfer (93.33\% precision across NLP/vision/training domains) reveals fundamental deep learning design confounds generalize beyond modality boundaries. Automated knowledge base construction achieves 84\% coverage without manual curation.

By treating testability as a verifiable predicate---∃ (D,B,M) in knowledge base---rather than expert intuition, this work demonstrates proof-of-concept for meta-research tool evaluation with the same rigor as primary research. Future work will close the external validity gap through real-world experiment execution, extend KB coverage to 95\%+ via multi-source aggregation, and improve recall from 10\% to 70-80\% through semantic extraction while maintaining low false positive rates.
